\section{Conclusions}
\label{sec:conclusion}

Class imbalance and class overlap are not interchangeable problems. In the synthetic-to-real factorial experiment, explicit separability pressure explains the large Vigitel-obesity utility change but not the COVID-19 variation. In four-task augmentation, SepAware improves CTGAN for COVID-19 and both Vigitel outcomes but not kidney mortality, is inconsistent for TVAE, and is often matched or exceeded by simple rebalancing. Closer minority coverage is accompanied by reduced diversity.

SepAware should therefore be understood as an auditable mechanism for controlling and studying class-conditional geometry, not as a universally superior augmentation method. Its practical use requires evidence that a predictive gain survives simple baselines and does not come from removing difficult but legitimate minority cases. The four-task study establishes this conditional direction and identifies kidney mortality as a clear nonbeneficial boundary case.
